% 两级 Harness 数据流图（详细版，含全部字段）
% 编译：xelatex -interaction=nonstopmode dataflow.tex
% 渲染：pdftoppm -png -r 200 dataflow.pdf dataflow
\documentclass[border=8pt]{standalone}
\usepackage{ctex}
\usepackage{tikz}
\usetikzlibrary{positioning, arrows.meta, calc}

\definecolor{runbg}{HTML}{FFF8E1}
\definecolor{epbg}{HTML}{E3F2FD}
\definecolor{ch1}{HTML}{C62828}
\definecolor{ch2}{HTML}{2E7D32}
\definecolor{ch3}{HTML}{6A1B9A}
\definecolor{statelabel}{HTML}{37474F}

\tikzset{
  nodebox/.style={draw, rounded corners=2pt, fill=white, font=\small,
                  text width=2.4cm, align=center, inner sep=3pt},
  statebox/.style={draw, rounded corners=3pt, fill=white, font=\scriptsize,
                   text width=4.8cm, align=left, inner sep=4pt},
  basebox/.style={draw, rounded corners=2pt, fill=white, font=\small,
                  text width=3.0cm, align=center, inner sep=4pt},
  arrow/.style={-{Stealth[length=2.5mm]}, thick},
  charrow/.style={-{Stealth[length=2.5mm]}, line width=1.3pt},
  dashedarrow/.style={-{Stealth[length=2.5mm]}, dashed, line width=1pt},
}

\begin{document}
\begin{tikzpicture}

% ============ 背景区域（直接矩形，可控） ============
\fill[runbg, opacity=0.55, rounded corners=4pt]
  (-0.2,0.1) rectangle (23.8,-5.4);
\fill[epbg, opacity=0.55, rounded corners=4pt]
  (-0.2,-5.6) rectangle (23.8,-12.2);

\draw[line width=1.2pt, rounded corners=4pt, dashed, gray!60]
  (-0.2,0.1) rectangle (23.8,-5.4);
\draw[line width=1.2pt, rounded corners=4pt, dashed, gray!60]
  (-0.2,-5.6) rectangle (23.8,-12.2);

% ============ Run 区标签 ============
\node[font=\large\bfseries, anchor=south] (runtitle)
  at (11.8, 0.1) {RunHarness（主 agent）——一个 run = 完整一局 Pokemon 游戏};

% ============ RunState 字段盒 ============
\node[statebox, anchor=north west] (RS) at (0.1, -0.3) {
  {\bfseries RunState（大图状态）}\\[1mm]
  {\color{statelabel}\bfseries 身份（跨轮）}\\[-2pt]
  \texttt{run\_id}: str\\
  \texttt{goals}: list[TaskForHarness]\\
  \hspace*{2em}\textbf{栈顶 = goals[-1]}\\
  \texttt{outcomes}: list[HarnessEpisodeOutcomeResp]\\[1mm]
  {\color{statelabel}\bfseries 流转（单轮）}\\[-2pt]
  \texttt{episode\_id}: str $|$ None\\
  \texttt{outcome}: HarnessEpisodeOutcomeResp $|$ None\\[1mm]
  {\color{statelabel}\bfseries 结束}\\[-2pt]
  \texttt{done}: bool \quad \texttt{why}: str
};

% ============ Run 区四节点 + END ============
\node[nodebox] (begin) at (6.8, -2.0) {begin\\[-1pt] {\scriptsize 初始栈入 state}};
\node[nodebox] (plan) at (10.3, -2.0) {plan\\[-1pt] {\scriptsize 读 trace（mask）\\压栈/拆解\\[-1pt]{\color{gray}\scriptsize 静态：栈非空即派发}}};
\node[nodebox] (dispatch) at (14.0, -2.0) {dispatch\\[-1pt] {\scriptsize 生成 episode\_id\\栈顶交给子 agent\\[-1pt]{\color{ch1}\scriptsize AgentError→失败结算}}};
\node[nodebox] (reflect) at (17.7, -2.0) {reflect\\[-1pt] {\scriptsize 收结算→弹栈\\重压/结束}};
\node[nodebox, fill=gray!10] (endrun) at (21.2, -2.0) {END};

% 边
\draw[arrow] (begin) -- (plan);
\draw[arrow] (plan) -- (dispatch);
\draw[arrow] (dispatch) -- (reflect);
\draw[arrow] (reflect) -- (endrun);
% reflect→plan 回环
\draw[arrow] (reflect.south) -- ++(0,-0.7) -| (plan.south);
\node[font=\scriptsize, fill=white, inner sep=1pt]
  at ($(reflect.south)!0.5!(plan.south)+(0,0.5)$) {栈非空};
\node[font=\scriptsize, fill=white, inner sep=1pt]
  at ($(endrun.north)+(0,0.15)$) {栈空};

% ============ 通道①：调用（dispatch → episode） ============
\draw[charrow, color=ch1] (dispatch.south) -- ++(0,-0.7) -|
  node[pos=0.25, above, font=\scriptsize\color{ch1}, align=center]
  {\textbf{① 调用}\\\texttt{run(episode\_id, 栈顶, 全栈)}\\$\rightarrow$ \texttt{HarnessEpisodeOutcomeResp}}
  ($(look.north)+(-0.3,0.05)$) -- (look.north);

% ============ Episode 区标签 ============
\node[font=\large\bfseries, anchor=south] (eptitle)
  at (11.8, -5.6) {EpisodeHarness（子 agent）——一个 episode = 解决栈顶一个目标};

% ============ EpisodeRunState 字段盒 ============
\node[statebox, anchor=north west] (EPS) at (0.1, -5.9) {
  {\bfseries EpisodeRunState（小图状态）}\\[1mm]
  {\color{statelabel}\bfseries 身份（跨步）}\\[-2pt]
  \texttt{episode\_id}: str\\
  \texttt{task}: TaskForHarness\\
  \texttt{step}: int \quad \textbf{全项目唯一}\\
  \texttt{goals}: list[Goal] \quad 只读投影\\[1mm]
  {\color{statelabel}\bfseries 流转（单步）}\\[-2pt]
  \texttt{observation}: ObservationFromWorld\\
  \texttt{action\_space}: ActionSpaceForBrain\\
  \texttt{memories}: list[StepMemory]\\
  \texttt{action}: Action\\
  \texttt{pending\_observation}: ObservationFromWorld
};

% ============ Episode 区六节点 + END ============
\node[nodebox] (look) at (6.8, -8.0) {look\\[-1pt]{\scriptsize 接帧·盖章\\判成败→OBSERVE}};
\node[nodebox] (retrieve) at (10.3, -8.0) {retrieve\_memory\\[-1pt]{\scriptsize 四类查询\\本局/object/知识/摘要}};
\node[nodebox] (think) at (13.8, -8.0) {think\\[-1pt]{\scriptsize choose()\\$\rightarrow$Action}};
\node[nodebox] (press) at (17.3, -8.0) {press\\[-1pt]{\scriptsize execute()\\推进世界→新观测}};
\node[nodebox] (summarize) at (6.8, -10.5) {summarize\\[-1pt]{\scriptsize 蒸馏→摘要\\EPISODE\_MEMORY\_WRITE}};
\node[nodebox, fill=gray!10] (endep) at (10.3, -10.5) {END};
\node[nodebox] (remember) at (13.8, -10.5) {remember\\[-1pt]{\scriptsize 写 StepMemory\\STEP\_MEMORY\_WRITE}};

% 小图边
\draw[arrow] (look) -- (retrieve);
\draw[arrow] (retrieve) -- (think);
\draw[arrow] (think) -- (press);
\draw[arrow] (press.south) -- ++(0,-0.5) -| (remember.north);
\draw[arrow] (remember.west) -| (look.south);
\node[font=\scriptsize, fill=white, inner sep=1pt] at ($(press)!0.5!(remember)+(0,0.15)$) {else};

% look 的 done 分支
\draw[arrow] (look.south) -- (summarize.north);
\node[font=\scriptsize, fill=white, inner sep=1pt] at ($(look)!0.5!(summarize)+(-0.5,0)$) {done};
\draw[arrow] (summarize) -- (endep);

% ============ 通道②：Trace（右侧大弧 + 独立标注） ============
\draw[dashedarrow, color=ch2] (endrun.east) -- ++(0.5,0)
  -- ++(0,-5.3) -| (endep.east);
\node[font=\scriptsize\color{ch2}, align=left, text width=4.2cm]
  at (24.2, -4.0)
  {\textbf{② Trace（共享实例）}\\episode 写事件流\\run 按 \texttt{RUN\_TRACE\_MASK} 读\\\{START, END, ERROR, EPISODE\_MEMORY\_WRITE\}};

% ============ World / Memory 盒 ============
\node[basebox, text width=5cm] (world) at (2.8, -14.0) {
  {\bfseries World（world 实现层）}\\
  PyBoy 模拟器 + VLM 视觉\\
  \texttt{reset / step / observe}\\
  产出 \texttt{ObservationFromWorld}
};
\node[basebox, text width=6.5cm] (memory) at (10.0, -14.0) {
  {\bfseries Memory（共享实例）}\\
  三种记忆：\texttt{StepMemory / EpisodeMemory / ObjectMemory}\\
  落盘：\texttt{\{filename\}.md} = frontmatter 元数据 + 正文\\
  trace 另存完整副本；启动扫描目录重建索引
};

% ============ 通道③：Memory（共享） ============
\draw[charrow, color=ch3, dashed] (press.south) -- ++(0,-3.3) -| (world.north);
\draw[charrow, color=ch3, dashed] (remember.south) -- ++(0,-2.3) -| (memory.north);
\node[font=\scriptsize\color{ch3}, align=left, text width=4cm] at (18.5, -14.0)
  {\textbf{③ Memory 共享}\\episode 蒸馏写/检索读\\（跨 episode 连续性）};

\end{tikzpicture}
\end{document}
